\documentclass[exercice,a4]{thedocument}%
\usepackage[french]{babel}%
\usepackage{scratch3}%
\usepackage{tabularx}%
\usepackage{texmaths-symbols}
\startdatakeys
    \source{21GENMATAA2}
    \BO{2026}
    \cycle{4}
    \competencesDuSocle{CA,RA}
    \objectifsApprentissage{C4D2ba}
    \motsCles{Scratch,programme}
\enddatakeys
\begin{document}%

\medskip

On travaille avec le logiciel Scratch dont voici plusieurs copies d’écran :

\begin{center}
\begin{tabularx}{\linewidth}{|>{\footnotesize}X|>{\footnotesize}m{5cm}|}\hline
Script principal :&Bloc du motif A :\\
\begin{scratch}
\blockinit{Quand \greenflag est cliqué}
\blocklook{Initialisation}
\blockvariable{mettre \selectmenu{1\up{er} nombre} à \ovaloperator{nombre aléatoire} entre \ovalnum{1} et \ovalnum{10}}
\blockvariable{mettre \selectmenu{2\up{e} nombre} à \ovaloperator{nombre aléatoire} entre \ovalnum{1} et \ovalnum{10}}
\blockifelse{si \booloperator{\ovalmove{1\up{er} nombre} $<$ \ovalmove{2\up{e} nombre}} alors}
{\blocklook{Motif A}}
{\blocklook{Motif B}}
\end{scratch}
&\begin{scratch}
\initmoreblocks{définir \namemoreblocks{Motif A}}
\blockpen{effacer tout}
\blockpen{stylo en position d'écriture}
\blockmove{s’orienter à \ovalnum{90}}
\blockrepeat{répéter \ovalnum{2} fois}
{\blockmove{avancer de \ovalnum{30}}
\blockmove{tourner \turnleft{} de \ovalnum{90} degrés}
\blockmove{avancer de \ovalnum{30}}
\blockmove{tourner \turnright{} de \ovalnum{90} degrés}
}
\blockpen{relever le stylo}
\end{scratch}\\ \hline
\end{tabularx}
\end{center}

Le bloc \begin{scratch}\blocklook{Motif B} \end{scratch} permet de tracer un triangle.

\begin{center}
\begin{tabularx}{\linewidth}{|X|}\hline
\multicolumn{1}{|c|}{\textbf{Informations}}\\
\setscratch{scale=0.7}
\begin{itemize}
\item[$\bullet~~$]Le bloc \begin{scratch}\blocklook{Initialisation}\end{scratch} efface l’écran et prépare le lutin.
\item[$\bullet~~$]L’instruction \begin{scratch}\blockmove{s’orienter à  \ovalnum{90}}\end{scratch}
signifie que le lutin se dirige horizontalement vers la droite.
\item[$\bullet~~$]L’instruction \begin{scratch}\blockvariable{nombre aléatoire entre \ovalnum{1} et \ovalnum{10}}\end{scratch} permet de choisir au hasard un nombre
entier dans la liste suivante : 1 ; 2 ; 3 ; 4 ; 5 ; 6 ; 7 ; 8 ; 9 ; 10.
\item[$\bullet~~$]\og 1\up{er} nombre \fg{} et \og 2\up{e} nombre \fg{} sont deux variables qui s’affichent à l’écran.

À l’écran, l’affichage \begin{scratch}\blockvariable{1\up{er} nombre \ovalnum{4}}\end{scratch} indique que la variable \og 1\up{er} nombre \fg prend la valeur 4.
\end{itemize}\\ \hline
\end{tabularx}
\end{center}

\begin{enumerate}
\item Tracer à main levée une allure du motif A défini par le bloc \og Motif A \fg.
\item Après avoir cliqué sur le drapeau vert, l’écran affiche :

\setscratch{scale=0.7}\begin{scratch}\blockvariable{1\up{er} nombre \ovalnum{7}}\end{scratch}

\begin{scratch}\blockvariable{2\up{e} nombre \ovalnum{3}}\end{scratch}

Quel motif est alors affiché à l’écran : le \og Motif A \fg{} ou le \og Motif B \fg{} ?
\item On relance le programme.

Si la variable \og 1\up{er} nombre \fg{} prend la valeur 3, calculer la probabilité pour que l’écran affiche le \og Motif A \fg.
\end{enumerate}

\bigskip

\showthesource
\end{document}
